% ============================================================================
% Fichier  : partie6/P06_C20_incidents_progres.tex
% Titre    : Gérer les Incidents et Progresser
% Auteur   : MINKA MI NGUIDJOÏ Thierry Emmanuel
% ============================================================================
% Sources : ISO/IEC 27035 ; ISO/IEC 27001 clause 10. Mairie d'Eséka : fictif.
% ============================================================================

\chapter{Gérer les Incidents et Progresser}
\label{chap:incidents}
\chaptermark{Incidents et progrès}

\epigraph{%
    « Un incident bien géré n'est pas un échec~:\\
    c'est le moment où le système apprend. »%
}{Principe du cours}

\niveauChapitre{Intermédiaire}{La partie VI.}

\begin{boxobjectifs}
\begin{itemize}
    \item Décrire le processus de gestion des incidents.
    \item Classer et prioriser un incident.
    \item Conduire un retour d'expérience (\sigle{RETEX}).
    \item Relier l'incident à l'amélioration continue (clause~10).
    \item Mesurer les progrès du \sigle{SMSI} dans la durée.
\end{itemize}
\end{boxobjectifs}

\bigskip

% ============================================================================
\section*{Question directrice : comment transformer un incident en progrès ?}
% ============================================================================

L'histoire de la mairie a commencé par un incident. Ce dernier chapitre montre
comment un \sigle{SMSI} gère les incidents et s'en sert pour progresser,
bouclant le cycle du cours.

\medskip
\textbf{Plan du chapitre.}
\begin{enumerate}
    \item Le processus de gestion (section~\ref{sec:c20-processus}).
    \item Classer et prioriser (section~\ref{sec:c20-classer}).
    \item Le retour d'expérience (section~\ref{sec:c20-retex}).
    \item L'incident, moteur d'amélioration (section~\ref{sec:c20-amelioration}).
    \item Mesurer les progrès (section~\ref{sec:c20-progres}).
\end{enumerate}

\separateur

% ============================================================================
\section{Le Processus de Gestion}
\label{sec:c20-processus}
% ============================================================================

\begin{boxdef}
La \termecle{gestion des incidents}~: processus qui détecte, signale, qualifie,
traite et clôt les incidents de sécurité, en tirant les leçons de chacun
(\norme{ISO/IEC 27035}).
\end{boxdef}

\begin{tableaucours}{L{3.0cm}L{9.2cm}}{%
    \tableauHeader{Étape} &
    \tableauHeader{Contenu}
}
Détecter et signaler & Repérer et alerter, sans crainte de sanction \\
Qualifier & Confirmer l'incident, en mesurer la portée \\
Traiter & Contenir, éradiquer, rétablir \\
Clore et apprendre & Documenter et conduire le RETEX \\
\end{tableaucours}
\vspace{-0.8em}
\begin{center}{\small\itshape Tableau~20.1. Le processus de gestion des incidents}\end{center}

Le signalement est la clé~: un incident tu ne peut être traité. La culture du
signalement sans blâme, posée à la sensibilisation (Chapitre~\ref{chap:orga}),
porte ici ses fruits.

% ============================================================================
\section{Classer et Prioriser}
\label{sec:c20-classer}
% ============================================================================

\begin{boxprocedure}
\textbf{Constatation.} La mairie reçoit deux signalements~: un courriel
d'hameçonnage non ouvert, et un poste d'état civil qui redémarre seul.

\medskip
\textbf{Analyse.} On classe selon l'\emph{impact} et l'\emph{urgence}~: le
courriel non ouvert est mineur (surveiller)~; le poste d'état civil qui
redémarre touche un actif critique (priorité haute, isoler et investiguer).
Prioriser, c'est mettre les moyens là où l'enjeu est le plus fort.
\end{boxprocedure}

% ============================================================================
\section{Le Retour d'Expérience}
\label{sec:c20-retex}
% ============================================================================

\begin{boxdef}
\termecle{Retour d'expérience} (\sigle{RETEX})~: analyse, après un incident, de
ce qui s'est passé, de ce qui a bien ou mal fonctionné, et des améliorations à
apporter. Il vise la cause, non le coupable.
\end{boxdef}

Le \sigle{RETEX} de l'incident initial de la mairie a produit des mesures
durables~: interdiction des clés personnelles, sauvegardes 3-2-1, comptes
nominatifs, sensibilisation. L'incident, bien analysé, a été le point de départ
de tout le \sigle{SMSI}.

\begin{boxmetier}{RSSI}
\textbf{Chercher la cause, pas la faute.} La RSSI conduit le \sigle{RETEX} dans
un esprit d'amélioration~: on cherche pourquoi le système a permis l'incident,
non qui a cliqué. Un RETEX accusateur tarit les signalements~; un RETEX
constructif les encourage.
\end{boxmetier}

% ============================================================================
\section{L'Incident, Moteur d'Amélioration}
\label{sec:c20-amelioration}
% ============================================================================

L'incident alimente la clause~10 (amélioration)~: chaque \sigle{RETEX} génère
des actions correctives qui éliminent des causes, et parfois des améliorations
qui renforcent au-delà. C'est le « Act » du \sigle{PDCA}~: on referme la boucle
et l'on relance un cycle plus mûr.

% ============================================================================
\section{Mesurer les Progrès}
\label{sec:c20-progres}
% ============================================================================

\begin{boxobjet}{Le SMSI de la mairie, un an après}
\textbf{Constatation.} Un an après l'incident, la mairie a~: une politique
signée, des risques appréciés, des mesures déployées, un tableau de bord, un
audit interne passé, une certification en vue.

\medskip
\textbf{Analyse.} Le progrès se mesure~: délai de reprise réduit, sauvegardes
testées, agents formés, incidents mieux traités. La mairie ne promet pas le
risque zéro~; elle démontre un risque \emph{maîtrisé}. C'est l'aboutissement du
cours~: non pas un système parfait, mais un système qui apprend.
\end{boxobjet}

\separateur

\begin{boxresume}
\textbf{Ce qu'il faut retenir de ce chapitre~:}
\begin{itemize}
    \item La \textbf{gestion des incidents} détecte, signale, qualifie, traite
          et clôt.
    \item On \textbf{classe} et \textbf{priorise} selon l'impact et l'urgence.
    \item Le \textbf{RETEX} cherche la cause, non le coupable.
    \item L'incident nourrit l'\textbf{amélioration continue} (clause~10).
    \item Le progrès se \textbf{mesure}~: un \sigle{SMSI} mûr maîtrise le
          risque, sans l'annuler.
\end{itemize}
\end{boxresume}

\begin{boxinfo}
\textbf{Points de contrôle de l'élève.} Je sais~: décrire le processus de
gestion des incidents~; classer et prioriser~; conduire un RETEX~; relier
l'incident à l'amélioration~; mesurer les progrès du \sigle{SMSI}.
\end{boxinfo}

\bigskip

\begin{boxexo}[title={\ding{45}\quad Exercice~20.1. Le processus \niveaubase}]
Énoncez les quatre étapes de la gestion des incidents.
\end{boxexo}

\begin{boxexo}[title={\ding{45}\quad Exercice~20.2. Signaler \niveaubase}]
Pourquoi le signalement sans blâme est-il essentiel~?
\end{boxexo}

\begin{boxexo}[title={\ding{45}\quad Exercice~20.3. RETEX \niveaubase}]
Qu'est-ce qu'un retour d'expérience, et que vise-t-il~?
\end{boxexo}

\begin{boxexo}[title={\ding{45}\quad Exercice~20.4. Prioriser \niveaumoyen}]
Entre le courriel non ouvert et le poste qui redémarre, lequel traiter d'abord~?
Pourquoi~?
\end{boxexo}

\begin{boxexo}[title={\ding{45}\quad Exercice~20.5. Cause ou faute \niveaumoyen}]
Pourquoi un RETEX accusateur tarit-il les signalements~?
\end{boxexo}

\begin{boxexo}[title={\ding{45}\quad Exercice~20.6. Incident et clause 10 \niveaumoyen}]
Comment un incident alimente-t-il l'amélioration continue~?
\end{boxexo}

\begin{boxexo}[title={\ding{45}\quad Exercice~20.7. L'incident fondateur \niveaumoyen}]
Montrez comment l'incident initial de la mairie a produit des mesures durables.
\end{boxexo}

\begin{boxexo}[title={\ding{45}\quad Exercice~20.8. Mesurer le progrès \niveaumoyen}]
Citez trois indicateurs montrant que le \sigle{SMSI} de la mairie a progressé.
\end{boxexo}

\begin{boxexo}[title={\ding{45}\quad Exercice~20.9. Procédure d'incident \niveauexpert}]
Rédigez la procédure de gestion des incidents de la mairie d'Eséka (détection à
clôture).
\end{boxexo}

\begin{boxexo}[title={\ding{45}\quad Exercice~20.10. RETEX complet \niveauexpert}]
Conduisez le RETEX écrit de l'incident initial de la mairie (déroulé, causes,
actions correctives, améliorations).
\end{boxexo}

\bigskip

\begin{boxinfo}
\textbf{Pour aller plus loin.}
\begin{itemize}
    \item ISO/IEC 27001, clause~10~\cite{iso27001}~; auditologie~\cite{minka2026auditologie}.
\end{itemize}
\end{boxinfo}

\begin{boxinfo}
\textbf{Clôture du cours.} De l'incident à la certification, la mairie d'Eséka
a appris à \emph{décider} sa sécurité, à la \emph{tenir} dans la durée et à la
\emph{prouver}. Tel est le management de la sécurité de l'information~: non pas
un état que l'on atteint, mais une roue que l'on fait tourner. Planifier,
déployer, contrôler, agir~: et recommencer, un cran plus haut.
\end{boxinfo}
